% =====================================================================
%  Chapter 18 — The Normal-Inverse-Gamma Posterior in Detail
% =====================================================================

\chapter{The Normal-Inverse-Gamma Posterior in Detail}\idx{Normal-inverse-gamma}
\label{ch:nig}

Chapter 17 stated the model and its update. This chapter takes the two
scalar parameters seriously, because the book's central claim --- that
the machine's job is to estimate the posterior over the edge --- lives
entirely in the relationship between \(\alpha_n\), \(\beta_n\) and the
quantities the portfolio layer consumes.

The two numbers that govern the prior are hyperparameters\idx{Hyperparameters}, and they are
not themselves estimated from the data: they are set before the data
arrives, and they are what the trial family varies. \(\alpha_0\) and
\(\beta_0\) set the prior on the scale, and \(\lambda\) sets the prior
precision on each coefficient.

\section{Why inverse-gamma}

The scale \(\sigma^2\) is unknown and the coefficients are unknown, and
they are not independent: a large \(\sigma^2\) makes any particular
coefficient hard to pin down. The conjugate treatment is to give
\(\sigma^2\) an inverse-gamma prior and condition the coefficient prior
on it, so that \(\beta \mid \sigma^2\) is normal with precision proportional
to \(1/\sigma^2\).

With \(\alpha_0 = 2.5\) and \(\beta_0 = 0.5\), the prior mean of
\(\sigma^2\) is
\begin{equation}
E[\sigma^2] = \frac{\beta_0}{\alpha_0 - 1} = \frac{0.5}{1.5} = \frac{1}{3},
\end{equation}
which in monthly terms is an annualised volatility of about
\(\sqrt{12/3} = 2\), i.e. 200 per cent. That is a deliberately absurd
prior on the scale: it says that before seeing any data, the residual
volatility of a fund could be anything, and that the data must supply
almost all of the information.

The reason for the absurdity is that the prior on the scale is the one
place where a sensible prior would be actively harmful. A prior that said
"monthly residual volatility is around four per cent" would dominate the
data in the early years, when there are few observations, and would
suppress exactly the volatility estimate the risk layer needs. An
uninformative prior on the scale costs almost nothing, because the data
speaks quickly about the scale --- the scale needs only the number of
observations, not their signal.

\section{The posterior scale and what it measures}

The posterior parameters are

\begin{equation}
\alpha_n = \alpha_0 + \frac{n}{2},
\qquad
\beta_n = \beta_0 + \frac{1}{2}\,\mathrm{RSS}_n,
\end{equation}
where \(\mathrm{RSS}_n\) is the residual sum of squares plus the prior
term, and the posterior mean of \(\sigma^2\) is

\begin{equation}
E[\sigma^2 \mid y] = \frac{\beta_n}{\alpha_n - 1}.
\end{equation}

What this quantity is \emph{not} is the volatility of the fund. It is the
volatility of the residual of the regression --- the part of the fund's
monthly return that the twelve features do not explain. For a broad
market tracker, the features explain almost nothing, so the residual
volatility is nearly the whole volatility. For a fund whose recent
behaviour is well described by its own momentum and volatility, the
residual is smaller.

This matters because the book's risk layer needs the \emph{total} risk of
the portfolio, and gets it from the empirical covariance matrix of
returns, not from the model's residual scale. The model's scale is used
for the thing it is good at: measuring how much of the forecast's
uncertainty is irreducible.

\section{Where the doubt comes from}

The posterior variance of the mean at a new point is
\begin{equation}
v_{\mathrm{mean}} = \hat{\sigma}^2\, x_*^{\prime}A_n^{-1}x_*.
\end{equation}

Added to it is the residual scale, and the sum is the variance of the
\emph{posterior predictive\idx{Posterior predictive}} distribution for a fund the model has not seen
before. The forecast therefore arrives with a mean and a variance, and the
portfolio layer consumes only those two moments; it never sees the
distribution itself.

Three things determine its size, and they are worth separating because
they behave differently over the life of the system.

\textbf{The information in the features.} The term \(x_*^{\prime}A_n^{-1}x_*\)
is the Mahalanobis distance of the new observation from the prior mean in
the metric of the data's information. A fund whose features are unusual
relative to everything the model has seen --- a fund that has just moved
far from its own history --- has a larger value, and therefore a more
uncertain forecast.

\textbf{The amount of data.} Through the prior term in
\(A_n^{-1} = (\Lambda + X^{\prime}X)^{-1}\), the prior acts as a floor. In
the early years, when \(X^{\prime}X\) is small, \(A_n^{-1}\) is close to
\(\Lambda^{-1}\), and every forecast carries the prior's uncertainty
however much data there is. As the panel grows, \(A_n^{-1}\) shrinks and
the forecasts tighten --- but only for the funds whose features are
well-represented in the data.

\textbf{Disagreement with the prior.} The prior mean is zero. A forecast
far from zero is a forecast the prior does not support, and the model's
own uncertainty about it is correspondingly larger.

The posterior band and the share of the uncertainty that is doubt about the forecast are shown in \figref{fig:posterior}.

\section{The consequence that shapes the portfolio layer}

Here is the step that the rest of the book depends on, and it is worth
stating as a single sentence.

The forecast \(\mu\) is a random variable whose posterior spread is
\(v_{\mathrm{mean}}\). The investor does not know \(\mu\); the investor
knows a distribution. If the portfolio layer sized a position on \(\mu\)
alone, it would be acting as though the investor knew the expected return
exactly, and would therefore take a full position whenever the sign of
\(\mu\) was right --- which is precisely when it is most likely to be
wrong.

The certainty-equivalent construction of Chapter 26 does the opposite: it
replaces the unknown expected return by the quantity an investor who is
risk-averse to being wrong should act on, and the size of the position
falls as the doubt rises. When the doubt is large --- which, as
\figref{fig:posterior} shows, it usually is --- the position
becomes small, and the portfolio is close to the stated policy.

This is why the book can report, in Part VI, that the machine's
contribution to the return is small, and still be describing a system that
works. A layer that contributes nothing when it knows nothing is not a
defect; it is the layer doing its job.

\begin{keytakeaways}
\item The scale prior is deliberately uninformative, because the data
speaks quickly about the scale and a sensible prior would suppress the
early-year volatility estimates the risk layer needs.
\item The posterior scale is the volatility of the regression residual,
not the volatility of the fund; the portfolio's total risk comes from the
empirical covariance of returns.
\item The posterior variance of the mean is small relative to the
posterior variance of the observation, and that gap is the entire
justification for the certainty-equivalent sizing of Chapter 26.
\item A layer that contributes nothing when it knows nothing is the layer
working correctly.
\end{keytakeaways}
